\subsubsection*{6.3 Decodierung eines modifizierten ISBN-Codes}


Der ISBN-Code ist ein modulo-11 Code $\subseteq F_{11}^{\,10}$. Seine Kontrollmatrix besteht nur aus der Zeile
$\mathbf{h} = (1\ 2\ 3\ \dots\ 10)$, die letzte Ziffer $c_{10}$ eines Codewortes $\mathbf{c}$ ist die Prüfziffer.
Um einen Fehler zu korrigieren, muss die Fehlerstelle $x$, $1 \le x \le 10$, und die Fehlergröße $e_x$ berechnet werden.
Dazu sind zwei Gleichungen bzw.\ zwei Prüfziffern erforderlich. Setzen
\[
  H = \begin{pmatrix} 1 & 1 & 1 & \dots & 1\\ 1 & 2 & 3 & \dots & 10 \end{pmatrix},
\]
Formen $H$ in systematische Form $(A \mid E)$ um (Einheitsmatrix rechts), um die Generatormatrix zu bestimmen.
\[
  \to\qquad (A \mid E) = \begin{pmatrix}
    9 & 8 & 7 & 6 & 5 & 4 & 3 & 2 & 1 & 0\\
    3 & 4 & 5 & 6 & 7 & 8 & 9 & 10 & 0 & 1
  \end{pmatrix},
\]
% Leere Einträge der Matrix wie im Original (entsprechen 0).
\[
  \Rightarrow\qquad G = (E \mid -A^T) = \begin{pmatrix}
    1 & & & & & & & & 2 & 8\\
    & 1 & & & & & & & 3 & 7\\
    & & 1 & & & & & & 4 & 6\\
    & & & 1 & & & & & 5 & 5\\
    & & & & 1 & & & & 6 & 4\\
    & & & & & 1 & & & 7 & 3\\
    & & & & & & 1 & & 8 & 2\\
    & & & & & & & 1 & 9 & 1
  \end{pmatrix}
\]
\begin{align*}
  \Rightarrow\qquad & \mathbf{c} = \mathbf{a}\,G, \quad \text{also } c_i = a_i \text{ für } 1 \le i \le 8.\\
  & c_9 \equiv \sum_{i=1}^{8} (1+i)\,c_i \pmod{11}, \qquad c_{10} \equiv \sum_{i=1}^{8} (9-i)\,c_i \pmod{11}.
\end{align*}

Sei $\mathbf{r}$ wie in Abschnitt 1.3 das empfangene Wort. Falls $\mathbf{r} = \mathbf{c}$, dann gilt
% KORRIGIERT: "s(r)" statt einheitlich S(r)
\[
  S(\mathbf{r}) = H\,\mathbf{r}^T = H\,\mathbf{c}^T = \mathbf{o}
\]
Falls ein Fehler an der Fehlerstelle $x$ mit der Größe $e_x$ auftrat, sieht $\mathbf{r}$ so aus
\begin{align*}
  & \mathbf{r} = \mathbf{c} + \mathbf{e} \quad\text{mit}\quad \mathbf{e} = (0\ \dots\ 0\ e_x\ 0\ \dots\ 0)\\
  \Rightarrow\qquad & H\,\mathbf{r}^T = H\,\mathbf{e}^T = \begin{pmatrix} e_x\\ x e_x \end{pmatrix}
    = \begin{pmatrix} 1\\ x \end{pmatrix} e_x = \begin{pmatrix} s_1\\ s_2 \end{pmatrix}\\
  % KORRIGIERT: "x = s_2/s_1, c_x = r_x - e_x" ohne Modul und ohne Begründung für die Division
  & s_1 = e_x,\quad s_2 \equiv x \cdot e_x,\quad x \equiv s_2 \cdot s_1^{-1},\quad c_x \equiv r_x - e_x \pmod{11}
\end{align*}
($s_1 = e_x \neq 0$, da ein Fehler auftrat; wegen $\ggT(s_1, 11) = 1$ ist $s_1$ modulo $11$ invertierbar.)\\
d.h.\ $s_1$ ist die Fehlergröße und die Fehlerstelle ist $x \equiv s_2 \cdot s_1^{-1} \pmod{11}$.

\begin{beispiel}
Sei $\mathbf{a} = 0\ 0\ \dots\ 0$, 8 Nullen, dann ist $\mathbf{c} = 0\ 0\ \dots\ 0\ 0\ 0$, 10 Nullen.\\
Angenommen
\[ \mathbf{r} = 0\ 0\ \dots\ 0\ 5\ 0. \]
% KORRIGIERT: "(1 9)^T 5" ohne Malpunkt
Dann ist $S(\mathbf{r}) = H\,\mathbf{r}^T = \begin{pmatrix} 5\\ 45 \end{pmatrix} = \begin{pmatrix} 1\\ 9 \end{pmatrix} \cdot 5$, also $e_x = 5$, $x = 9$.\\
Tatsächlich ist $\mathbf{c} = \mathbf{o}$ Codewort in jedem linearen Code.
\end{beispiel}

\begin{beispiel}
Sei $\mathbf{r} = 3\ 5\ \dots\ 0\ 0\ 0$.\\
% KORRIGIERT: "=" statt "\equiv (mod 11)"; Kürzen von 2/8 und Inverse ohne Begründung
Dann ist $\mathbf{s} = H\,\mathbf{r}^T = \begin{pmatrix} 8\\ 13 \end{pmatrix} \equiv \begin{pmatrix} 8\\ 2 \end{pmatrix} \pmod{11}$, also $e_x = 8$,
$x \equiv 2 \cdot 8^{-1} \equiv 1 \cdot 4^{-1} \equiv 3 \pmod{11}$
(Kürzen mit $2$ erlaubt, da $\ggT(2, 11) = 1$; $4^{-1}$ existiert, da $\ggT(4, 11) = 1$).\\
Wegen $3 \cdot 4 = 12 \equiv 1 \pmod{11}$, also $c_3 \equiv 0 - 8 \equiv 3 \pmod{11}$
\[ \mathbf{c} = 3\ 5\ 3\ \dots\ 0\ 0\ 0. \]
% Im Original als "1.3 + 2.5 + 3.3" (Punkt als Malzeichen) geschrieben.
% KORRIGIERT: "(11 22)^T = (0 0)^T" statt \equiv (mod 11)
Tatsächlich ist $H\,\mathbf{c}^T = \begin{pmatrix} 3+5+3\\ 1\cdot 3 + 2\cdot 5 + 3\cdot 3 \end{pmatrix} = \begin{pmatrix} 11\\ 22 \end{pmatrix} \equiv \begin{pmatrix} 0\\ 0 \end{pmatrix} \pmod{11}$, also $\mathbf{c}$ ein Codewort.
\end{beispiel}

\begin{bemerkung}
% KORRIGIERT: "r" statt fett gedrucktem Vektor
Falls $s_1 = 0$ und $s_2 \neq 0$, enthält $\mathbf{r}$ mehr als einen Fehler.
\end{bemerkung}
